% Options for packages loaded elsewhere
\PassOptionsToPackage{unicode}{hyperref}
\PassOptionsToPackage{hyphens}{url}
\documentclass[
  english,
  11pt,
  a4paper,
]{article}
\usepackage{xcolor}
\usepackage[margin=23mm]{geometry}
\usepackage{amsmath,amssymb}
\setcounter{secnumdepth}{-\maxdimen} % remove section numbering
\usepackage{iftex}
\ifPDFTeX
  \usepackage[T1]{fontenc}
  \usepackage[utf8]{inputenc}
  \usepackage{textcomp} % provide euro and other symbols
\else % if luatex or xetex
  \usepackage{unicode-math} % this also loads fontspec
  \defaultfontfeatures{Scale=MatchLowercase}
  \defaultfontfeatures[\rmfamily]{Ligatures=TeX,Scale=1}
\fi
\usepackage{lmodern}
\ifPDFTeX\else
  % xetex/luatex font selection
  \setmainfont[]{DejaVu Serif}
  \setsansfont[]{DejaVu Sans}
  \setmonofont[Scale=0.82]{DejaVu Sans Mono}
\fi
% Use upquote if available, for straight quotes in verbatim environments
\IfFileExists{upquote.sty}{\usepackage{upquote}}{}
\IfFileExists{microtype.sty}{% use microtype if available
  \usepackage[]{microtype}
  \UseMicrotypeSet[protrusion]{basicmath} % disable protrusion for tt fonts
}{}
\makeatletter
\@ifundefined{KOMAClassName}{% if non-KOMA class
  \IfFileExists{parskip.sty}{%
    \usepackage{parskip}
  }{% else
    \setlength{\parindent}{0pt}
    \setlength{\parskip}{6pt plus 2pt minus 1pt}}
}{% if KOMA class
  \KOMAoptions{parskip=half}}
\makeatother
\usepackage{longtable,booktabs,array}
\usepackage{caption}
\captionsetup[table]{skip=6pt}
\newcounter{none} % for unnumbered tables
\usepackage{calc} % for calculating minipage widths
% Correct order of tables after \paragraph or \subparagraph
\usepackage{etoolbox}
\makeatletter
\patchcmd\longtable{\par}{\if@noskipsec\mbox{}\fi\par}{}{}
\makeatother
% Allow footnotes in longtable head/foot
\IfFileExists{footnotehyper.sty}{\usepackage{footnotehyper}}{\usepackage{footnote}}
\makesavenoteenv{longtable}
\ifLuaTeX
\usepackage[bidi=basic,shorthands=off]{babel}
\else
\usepackage[bidi=default,shorthands=off]{babel}
\fi
\ifPDFTeX
\else
\babelfont{rm}[]{DejaVu Serif}
\fi
\ifLuaTeX
  \usepackage{selnolig} % disable illegal ligatures
\fi
\setlength{\emergencystretch}{3em} % prevent overfull lines
\providecommand{\tightlist}{%
  \setlength{\itemsep}{0pt}\setlength{\parskip}{0pt}}
\usepackage{fancyhdr}
\usepackage{titlesec}
\usepackage{needspace}
\definecolor{reviewblue}{HTML}{174E67}
\AtBeginDocument{\hypersetup{colorlinks=true,linkcolor=reviewblue,urlcolor=reviewblue,citecolor=reviewblue}}
\pagestyle{fancy}
\fancyhf{}
\fancyhead[L]{\small\sffamily FEILIAN technical assessment}
\fancyhead[R]{\small\sffamily Version 1.0.0}
\fancyfoot[L]{\footnotesize\sffamily Mounir IDRASSI}
\fancyfoot[R]{\footnotesize\sffamily\thepage}
\setlength{\headheight}{15pt}
\renewcommand{\headrulewidth}{0.3pt}
\renewcommand{\footrulewidth}{0pt}
\fancypagestyle{plain}{\fancyhf{}\fancyfoot[L]{\footnotesize\sffamily Mounir IDRASSI}\fancyfoot[R]{\footnotesize\sffamily\thepage}\renewcommand{\headrulewidth}{0pt}}
\titleformat{\section}{\large\sffamily\bfseries\color{reviewblue}}{}{0pt}{}
\titlespacing*{\section}{0pt}{2.1ex plus 1ex minus .2ex}{1ex}
\AtBeginEnvironment{longtable}{\small}
\setlength{\emergencystretch}{3em}
\clubpenalty=10000
\widowpenalty=10000
\displaywidowpenalty=10000
\brokenpenalty=10000
\predisplaypenalty=10000
\usepackage{bookmark}
\IfFileExists{xurl.sty}{\usepackage{xurl}}{} % add URL line breaks if available
\urlstyle{same}
% fallback for those not using the hyperref driver hyperxmp:
\makeatletter
\@ifundefined{xmpquote}{\newcommand{\xmpquote}[1]{#1}}{}
\makeatother
\hypersetup{
  pdftitle={FEILIAN: technical assessment of the submitted specification and implementations},
  pdfauthor={Mounir IDRASSI (mounir@amcrypto.jp)},
  pdflang={en},
  hidelinks,
  pdfcreator={LaTeX via pandoc}}

\title{FEILIAN: technical assessment of the submitted specification and
implementations}
\usepackage{etoolbox}
\makeatletter
\providecommand{\subtitle}[1]{% add subtitle to \maketitle
  \apptocmd{\@title}{\par {\large #1 \par}}{}{}
}
\makeatother
\subtitle{Cryptographic review · Version 1.0.0}
\author{Mounir IDRASSI
(\href{mailto:mounir@amcrypto.jp}{\nolinkurl{mounir@amcrypto.jp}})}
\date{24 September 2026}

\begin{document}
\maketitle

\textbf{Abstract.} This review examines the FEILIAN hash submission
dated 30 June 2026, including its specification, six C implementations,
four RTL architectures and validation material. Static analysis
identifies a message-length binding defect in the 1SC, 4SC and 8SC
hardware wrappers; the 2SC counter path differs. The specification and
implementations disagree on security-relevant details. All six Appendix
B examples are exactly reproduced: the 512/768 examples use cumulative
padded-block lengths, while the 1024 examples use actual message
lengths. Verification covers all 45 printed round states and nine
initial states. Ordinary tests pass all 24,582 C-to-KAT comparisons.
Exact arithmetic confirms the Sigma maps' invertibility and intended
constants. Shared test-generation code also has portability and
error-propagation defects, distinct from the hash primitive. The
security discussion does not establish the claimed composition, margin
or application guarantees. No full-round cryptanalytic break of the
correctly encoded software hash is established. A corrected, unambiguous
submission and renewed independent review are recommended.

\textbf{Prior work and preparation.} Markku-Juhani O. Saarinen's
\href{https://ngcc.dev/reports/hash-10.html}{FEILIAN report of 23
September 2026}, credited there with AI assistance, was consulted before
this review. It reports silent allocation failure in the reference C
implementations. R5 extends that source-level observation to the
optimized implementations; it is not presented as an independently
discovered original finding. R1--R4 and R6--R8 describe additional
observations relative to that report as consulted on 24 September; no
exhaustive novelty claim is made. A second researcher subsequently
supplied a private review. That researcher is credited for the Appendix
B counter explanation and for raising the additional hardware and
shared-DRNG issues incorporated after critical checking. The second
review is not endorsed as a whole; its broader experimental claims are
not included in this report's verified coverage. The
\href{PROVENANCE.md}{provenance record} identifies the reviewed files
and the limits of verification.

\textbf{AI-use disclosure.} OpenAI Codex assisted with source and
specification analysis, mathematical checks, verification code and
publication preparation. The executable coverage and the static
conclusions are reported separately. This package is prepared for Mounir
IDRASSI's review and possible release; preparation does not assert that
he has independently verified every conclusion. See
\url{AI_DISCLOSURE.md}.

\clearpage

\section{Assessment and scope}\label{assessment-and-scope}

\textbf{Recommendation: withhold a positive security assessment and
integration approval for this submission.} There is a high-severity
message-length encoding defect in three supplied hardware
implementations, several mutually inconsistent definitions of the hash,
and substantial gaps in the argument supporting its security claims.
These are sufficient grounds to request a corrected submission and
renewed review. This review does \textbf{not} establish a full-round
cryptanalytic break of the hash defined by the written compression
algorithm or of the normally operating C implementations.

The most consequential concrete finding beyond the
\href{https://ngcc.dev/reports/hash-10.html}{existing allocation-failure
report} is at the boundary between zero padding, the bit counter, and
the RTL compression interface. The 1SC, 4SC, and 8SC cores omit the
current block's length from the compression input. The consequence is
loss of injective message encoding, independently of how strong the ARX
compression function is. The 2SC core contains a different, corrected
counter path.

The principal source is
\href{data/SOURCES.md\#source-1}{Specification.pdf}, dated 30 June 2026.
Page references below are the \textbf{printed page numbers}; the PDF
viewer's page number is one greater. Its SHA-256 is
\nolinkurl{8e3a9a109f3188ab57cbd0156c243f80a6453051cf98006dfef11d66bdeb5fa8}.
The \href{data/input_manifest.json}{input manifest} fingerprints all 74
original files. The archive digest on the public report was not checked
against an original ZIP, which was not present here.

\Needspace{7\baselineskip}

{\def\LTcaptype{none} % do not increment counter
\begin{longtable}[]{@{}
  >{\raggedright\arraybackslash}p{(\linewidth - 6\tabcolsep) * \real{0.0700}}
  >{\raggedright\arraybackslash}p{(\linewidth - 6\tabcolsep) * \real{0.3100}}
  >{\raggedright\arraybackslash}p{(\linewidth - 6\tabcolsep) * \real{0.3100}}
  >{\raggedright\arraybackslash}p{(\linewidth - 6\tabcolsep) * \real{0.3100}}@{}}
\toprule\noalign{}
\begin{minipage}[b]{\linewidth}\raggedright
ID
\end{minipage} & \begin{minipage}[b]{\linewidth}\raggedright
Finding
\end{minipage} & \begin{minipage}[b]{\linewidth}\raggedright
Assessment
\end{minipage} & \begin{minipage}[b]{\linewidth}\raggedright
Evidence
\end{minipage} \\
\midrule\noalign{}
\endhead
\bottomrule\noalign{}
\endlastfoot
R1 & Current-block length omitted by 1SC/4SC/8SC RTL & High severity for
those hardware hash implementations & Static dataflow and state-machine
analysis; no RTL simulation \\
R2 & Chapter 2, C, and RTL do not define one consistent algorithm &
Blocks conformance and transfer of security conclusions & Source
comparison, PDF rendering, independent model, exact constant checks \\
R3 & Appendix B uses padded counts for 512/768, true counts for 1024 &
Security-relevant encoding inconsistency; C/KATs use true counts & All
six digests, 45 round snapshots, nine initial states and nine domains
verified \\
R4 & C bit-input and length-conversion contracts are incomplete &
Correctness/security defect; platform-dependent severity & Static
analysis of all six C implementations \\
R5 & Allocation failure is silently accepted in optimized C too &
Extends the scope of the existing implementation finding & Static
analysis; allocation failure not induced \\
R6 & Proof and security-margin claims are not established & Major
assurance gap, not a demonstrated hash break & Mathematical and
compositional review \\
R7 & MAC/KDF/XOF recommendations exceed what was established &
Application specifications need separate review & Domain, entropy,
interface, and proof-scope analysis \\
R8 & Shared DRNG and KAT driver have portability and error-handling
defects & Test-generation portability and error handling & Static source
review and identical-file hashes; no failure injection \\
\end{longtable}
}

\Needspace{5\baselineskip}

\section{R1. Hardware message-length
binding}\label{r1.-hardware-message-length-binding}

Section 2.3, pp.~12--13, updates \texttt{t} to include the current block
\textbf{before} calling the compression function. This is a security
requirement because the message is padded with zeros and has no in-block
delimiter or encoded length.

In the \href{data/SOURCES.md\#source-2}{8SC core}, \texttt{cf\_domain}
is made from \texttt{hashed\_bits\_reg}. That register represents only
the completed previous blocks. Its increment occurs in the
\texttt{WAIT\_CF} completion branch at line 160, after the compression
output has already been calculated. The last-block flag is provided
correctly, but the current final-block length is absent from the domain
used to obtain the final digest. The compression module latches the
supplied domain on \texttt{start}, so a later counter update cannot
repair the result; see \href{data/SOURCES.md\#source-3}{compression.sv}.

The same defect occurs in the \href{data/SOURCES.md\#source-4}{1SC core}
and \href{data/SOURCES.md\#source-5}{4SC core}. Their domain assignments
are at line 135 and completion-time counter updates at line 157. Neither
interface supplies a separately authenticated final-block length to
compression.

Zero padding is not injective by itself. A final flag distinguishes the
role of a block, but does not distinguish all possible logical lengths
within that role. Removing the actual final length therefore permits
distinct logical inputs to have an identical encoded compression
transcript. No differential property, round reduction, birthday search,
or allocation failure is needed for this structural failure. It affects
variable-length hashing in these wrappers during normal execution. It
should be treated as a collision-resistance failure of their message
encoding, not as evidence that the specified ARX primitive has been
cryptanalytically broken. Protocols imposing an independently
authenticated, fixed input length would need a separate impact
assessment.

The \href{data/SOURCES.md\#source-6}{2SC core} is an important
exception: it computes \texttt{hashed\_bits\_updated} from the current
input length and passes that value into \texttt{make\_domain}. Do not
report this particular defect as affecting all four architectures.

The corrective requirement is to form one accepted block transaction
containing the block, final flag, and cumulative bit count
\textbf{including that block}, and preserve that tuple for the complete
compression operation. Counter overflow and invalid final byte counts
also require explicit rejection. The acceptance invariant is
\(t_{\mathrm{compression}}=t_{\mathrm{previous}}+\ell_{\mathrm{current}}\).
Checking only compression-round arithmetic will not detect an error in
the wrapper that supplies this count.

This finding follows from the source's dataflow. No RTL simulation was
performed in this review. The source-level conclusion is strong; a
corrected implementation still needs end-to-end RTL validation before
deployment.

\Needspace{5\baselineskip}

\section{R2. Inconsistent algorithm
definitions}\label{r2.-inconsistent-algorithm-definitions}

There are independent discrepancies, so correcting one typo is
insufficient.

\Needspace{7\baselineskip}

{\def\LTcaptype{none} % do not increment counter
\begin{longtable}[]{@{}
  >{\raggedright\arraybackslash}p{(\linewidth - 4\tabcolsep) * \real{0.2200}}
  >{\raggedright\arraybackslash}p{(\linewidth - 4\tabcolsep) * \real{0.3900}}
  >{\raggedright\arraybackslash}p{(\linewidth - 4\tabcolsep) * \real{0.3900}}@{}}
\toprule\noalign{}
\begin{minipage}[b]{\linewidth}\raggedright
Detail
\end{minipage} & \begin{minipage}[b]{\linewidth}\raggedright
Chapter 2 / specification
\end{minipage} & \begin{minipage}[b]{\linewidth}\raggedright
Supplied implementation evidence
\end{minipage} \\
\midrule\noalign{}
\endhead
\bottomrule\noalign{}
\endlastfoot
IV word 7 & \texttt{0x3D84D5B5B5470917}, p.~3 & C and RTL use
\texttt{0x3F84D5B5B5470917} \\
AddConstant placement & Tweak in row index 1, constants in row index 3,
p.~9 & All six C versions and all four RTL round implementations use
constants in row index 1 and tweak in row index 3 \\
Counter word order & \(t=t_0\mathbin{\|}t_1\); significance should be
made explicit alongside the little-endian rule & C and 2SC use high
word, then low word; 1SC/4SC/8SC use low word, then high word \\
Count supplied to compression & Includes the current block, p.~13 & C
and 2SC do; 1SC/4SC/8SC do not \\
Hardware variant labels & README labels 1SC and 4SC as 512-bit & Every
supplied \texttt{feilian\_pkg.sv} sets \texttt{DIGEST\_BITS\ =\ 1024}
and \texttt{VERSION\ =\ 0x400} \\
\end{longtable}
}

The AddConstant discrepancy was checked against a rendered PDF page, not
inferred solely from text extraction. The scalar implementation's
\href{data/SOURCES.md\#source-7}{round functions} are explicit. Appendix
A's C listing also follows the implementation's placement, so the
disagreement exists within the PDF itself.

The C \href{data/SOURCES.md\#source-8}{counter helper} stores the high
word in \texttt{hash\_bits{[}0{]}} despite later comments identifying
that slot as the low word. The SIMD \texttt{\_mm256\_set\_epi64x} call
agrees with the executed C ordering; some accompanying comments do not.
Correcting the RTL count timing without also settling word order will
leave a conformance failure.

Sage interval arithmetic confirmed that the C IV matches the first 256
hexadecimal fractional digits of π, including the \texttt{3F} word.
Exact integer square-root checks reproduced all four round constants.
This makes the printed \texttt{3D} value look like an erratum, but the
submitters must designate the canonical algorithm. The text's
description of these IV values as decimal digits is also inaccurate.

An independent Python implementation, using Chapter 2's IV and
AddConstant placement, little-endian message words from §3.2, and the C
high-word-first counter convention, produced different hashes from the C
code for the empty message, \texttt{abc}, and 129 repetitions of
\texttt{a}, for every digest size. This is a conformance comparison on
ordinary messages. For example, the first 16 digest bytes for
FEILIAN-1024(\texttt{abc}) are:

\Needspace{7\baselineskip}

{\def\LTcaptype{none} % do not increment counter
\begin{longtable}[]{@{}
  >{\raggedright\arraybackslash}p{(\linewidth - 2\tabcolsep) * \real{0.4300}}
  >{\raggedright\arraybackslash}p{(\linewidth - 2\tabcolsep) * \real{0.5700}}@{}}
\toprule\noalign{}
\begin{minipage}[b]{\linewidth}\raggedright
Interpretation
\end{minipage} & \begin{minipage}[b]{\linewidth}\raggedright
First 16 digest bytes, hexadecimal
\end{minipage} \\
\midrule\noalign{}
\endhead
\bottomrule\noalign{}
\endlastfoot
Supplied scalar/SIMD C & \nolinkurl{0507e4b1ef75d025dfa303ea0a8e2c47} \\
Chapter 2 under the stated decoding conventions &
\nolinkurl{8b27103d98eabd06b36a8a1d2148d8e6} \\
\end{longtable}
}

Security analysis must identify which version it covers. Tweak placement
and counter ordering affect the actual permutation family, including
which symmetry conditions correspond to valid message lengths. They
cannot be dismissed as output formatting differences.

The hardware digest formatting and MMIO word order need an explicit
byte-level contract as well. The cores reverse bytes within words and
reverse word positions when forming the packed digest, while the MMIO
port reads increasing packed slices. A driver may compensate, but no
such driver contract or end-to-end testbench is supplied.

The \href{data/SOURCES.md\#source-13}{top-level modules} declare
\texttt{DIGEST\_BITS} but instantiate \texttt{core\ u\_core} without
forwarding that parameter. The package's \texttt{VERSION} also remains
fixed, and the digest formatting and MMIO loops use the fixed
sixteen-word state width. Changing only the top-level digest width
therefore does not select FEILIAN-512 or FEILIAN-768 correctly. Width
propagation, version selection, digest selection and slice bounds must
be specified and reviewed together.

The hardware input interface supplies a byte count, with no partial-bit
count. Its supported domain is therefore byte-aligned messages; this is
an interface limitation, not evidence of mishandling a bit-length
argument that the hardware never accepts. The eight-bit
\texttt{valid\_bytes} field also lacks rejection of final-block values
greater than the 128-byte block capacity. For non-final blocks the
helper always adds 1,024 bits. The 2SC path uses the final count before
compression, whereas the other three defer it as described in R1. A
legal range, empty-message convention and rejection behavior must be
defined at the accepted-transaction boundary. These are static
observations; no out-of-range request was executed.

Reset and bus-handshake behavior also need an explicit integration
contract. A datapath register initialized on an accepted start before
producing valid output is not defective merely because it lacks reset. A
reset-required completion state or level-sensitive start is likewise
not, by itself, a security vulnerability. These observations do not add
separate reset or handshake findings.

\Needspace{5\baselineskip}

\section{R3. Appendix B uses two different counter
conventions}\label{r3.-appendix-b-uses-two-different-counter-conventions}

All six supplied C implementations passed every entry in their
corresponding \texttt{KAT\_2\_12} file: 4,097 lengths, from 0 through
4,096 bits. That is 24,582 successful C-to-KAT comparisons. The
independent Python model matched 20 selected KAT cases per version,
covering partial bytes and block boundaries, and all nine additional
ordinary example computations. Those C implementations use the
cumulative true message length.

The second reviewer identified an exact explanation for the four
previously unexplained Appendix B mismatches. Keeping the C IV, round
operations, version selection, high-word-first counter and digest
serialization, the following counts reproduce every printed example:

\Needspace{7\baselineskip}

{\def\LTcaptype{none} % do not increment counter
\begin{longtable}[]{@{}
  >{\raggedright\arraybackslash}p{(\linewidth - 6\tabcolsep) * \real{0.2200}}
  >{\raggedleft\arraybackslash}p{(\linewidth - 6\tabcolsep) * \real{0.1900}}
  >{\raggedright\arraybackslash}p{(\linewidth - 6\tabcolsep) * \real{0.2900}}
  >{\raggedright\arraybackslash}p{(\linewidth - 6\tabcolsep) * \real{0.3000}}@{}}
\toprule\noalign{}
\begin{minipage}[b]{\linewidth}\raggedright
Variant
\end{minipage} & \begin{minipage}[b]{\linewidth}\raggedleft
\texttt{abc}: final count (bits)
\end{minipage} & \begin{minipage}[b]{\linewidth}\raggedright
129 bytes of \texttt{a}: first / final count (bits)
\end{minipage} & \begin{minipage}[b]{\linewidth}\raggedright
Matching convention
\end{minipage} \\
\midrule\noalign{}
\endhead
\bottomrule\noalign{}
\endlastfoot
FEILIAN-1024 & 24 & 1,024 / 1,032 & Cumulative true message length \\
FEILIAN-768 & 1,024 & 1,024 / 2,048 & Cumulative padded-block length \\
FEILIAN-512 & 1,024 & 1,024 / 2,048 & Cumulative padded-block length \\
\end{longtable}
}

This was independently checked with the earlier review's model, against
a fresh extraction of the hash-identified PDF. All \textbf{six digests,
45 printed round snapshots and nine initial states} agree under the
assignments above. The check also recovers all nine domain values from
the first-round states and confirms their high-word-first counter
layout. It enforces the stated assignment for each variant, rather than
merely accepting whichever convention happens to match. The
\href{evidence/appendix_comparison.json}{complete evidence},
\href{data/appendix_states.json}{numeric transcriptions} and
\href{code/check_appendix.py}{bounded ordinary-example checker} are
included.

The four 512/768 digests still disagree with C. For example,
FEILIAN-512(\texttt{abc}) begins
\nolinkurl{b69fc86f312df402fd5d243c61ab3768} in Appendix B and
\nolinkurl{fc7320c565f60ec6a6ea26585f7c7d1f} in C. What is now
established is the precise numerical explanation. The historical
generator version and reason for this discrepancy remain unconfirmed
until the submitters provide provenance.

This is security-relevant encoding, not simply a different display
convention. Padded-block counts do not bind the actual final logical
length under zero-only padding. The final flag does not restore that
missing information. Any argument requiring injective message encoding
cannot be transferred to the convention illustrated by the 512/768
appendix examples. This observation concerns the appendix convention; it
does not show that the normally operating supplied C/KAT implementation
uses that convention, and it does not establish a full-round
cryptanalytic break of correctly encoded FEILIAN.

The initial IV matrices are displayed transposed relative to the numeric
row-major convention. The computed round states and the second-block
initial state are row-major. A uniform transpose of every displayed
matrix would be incorrect. The checker normalizes only the initial IV
display.

The correction must retain actual message-length binding and regenerate
all examples from the designated algorithm. Treating padded and true
counts as equally acceptable options would leave the security problem
unresolved. Agreement with KATs generated from the same implementation
demonstrates reproducibility, not independent agreement with the
specification or a valid security proof.

\Needspace{5\baselineskip}

\section{R4. Bit-input and length contracts in
C}\label{r4.-bit-input-and-length-contracts-in-c}

The same padding logic appears in all six C implementations. In the
scalar \href{data/SOURCES.md\#source-9}{padding function},
\texttt{memcpy} copies the complete final byte, but the unused bits of a
partial final byte are not masked. Those bits subsequently enter the
compression function. Either the API must expressly require a canonical
final byte and reject violations, or the implementation must mask the
bits outside \texttt{msg\_len\_bits}. The current header states only
that the argument is the number of message bits.

This is a failure to make the result depend only on the declared
bitstring. It is not evidence of a compression-function collision. The
KAT generator's DRNG masks unused bits itself, which explains why the
supplied partial-bit KATs do not test this API boundary; see
\href{data/SOURCES.md\#source-10}{drng.c}.

There is a separate allocation-size arithmetic issue. The code adds the
byte-rounded message length to the independently byte-rounded padding
length. For a non-byte-aligned message, that allocates one byte more
than the necessary whole-block buffer. This is harmless over-allocation
by itself; it should not be misreported as a proven out-of-bounds
access. The missing final-byte mask remains significant.

The \href{data/SOURCES.md\#source-11}{public API} accepts
\texttt{unsigned\ long\ long} bit lengths, then casts them to
\texttt{size\_t} without checking representability. On a platform with
narrower \texttt{size\_t}, the declared message length can be silently
truncated. This is a potentially serious correctness defect for a
supposedly portable implementation. On platforms where the widths match,
the subsequent \texttt{msg\_bits\ +\ 7} still needs an overflow check.
These conclusions are static; no 32-bit binary or extreme-length request
was run.

The presence of a two-word counter does not make this one-shot API
support the specification's entire domain of messages shorter than
\(2^{128}\) bits. It needs a streaming interface with validated counter
updates, or an expressly smaller supported-input limit. A streaming
implementation also avoids allocating a second copy of the complete
message.

Output-length semantics need clarification too. The instance fixes
\texttt{FEILIAN\_VERSION} at compile time but accepts any positive
\texttt{digest\_len\_bits} up to its maximum and truncates the digest. A
shortened result from the FEILIAN1024 implementation is therefore a
truncation of the 1024-version hash, not the independently versioned
FEILIAN512 hash. That can be a legitimate API choice, but it must not be
confused with selecting another family member. Reject unsupported named
lengths or expose raw truncation and family selection distinctly. For a
permitted non-byte-aligned output length, a buffer of
\(\lceil t/8\rceil\) bytes with unused bits cleared is a normal
representation of a \(t\)-bit digest; storage rounding alone is not a
vulnerability.

\Needspace{5\baselineskip}

\section{R5. Allocation-failure handling in optimized
C}\label{r5.-allocation-failure-handling-in-optimized-c}

Saarinen's \href{https://ngcc.dev/reports/hash-10.html}{23 September
report} identifies reference implementations returning successful
all-zero digests after padding allocation fails. Static inspection shows
the same path in the \href{data/SOURCES.md\#source-12}{SIMD
implementation}: the internal routine clears the digest and returns
without an error status, and the public wrapper returns success. The
512- and 768-bit optimized sources use the same logic.

The fix must propagate failure through the public return code and leave
the output unusable as a successful digest. Streaming reduces the
memory-pressure exposure but does not replace correct error propagation.
No allocation failure was induced during this review, and the public
report's witness was not rerun.

\Needspace{5\baselineskip}

\section{R6. Mathematical and security
arguments}\label{r6.-mathematical-and-security-arguments}

FEILIAN uses a 1,024-bit chaining state, a 1,024-bit message block, and
a public tweak containing version, final flag, and a 128-bit counter.
Its five phases each comprise one message-bearing round and three
message-free rounds. Each full round has two SubColumn/ShiftRow layers.
Compression applies a Davies--Meyer feedforward,
\(h' = E_{m,tw}(h) \oplus h\).

This architecture is intelligible and uses established ingredients. Its
security still depends on the concrete message- and tweak-dependent
permutation family. Naming HAIFA, Davies--Meyer, or Even--Mansour does
not establish that family behaves like the ideal object used in a
theorem.

First, §4.3.1, pp.~19--20, proves a statement about a compression
interface with a \textbf{salt} and final padding that encodes message
length. The preceding FEILIAN definition instead has a \textbf{final
flag}, no salt, and zero padding without an in-block length field. The
simulator validates the former padding and does not model the actual
flag. The proof therefore does not establish the stated result for the
specified construction. The original
\href{https://eprint.iacr.org/2007/278.pdf}{HAIFA paper} describes
explicit padding and salt conventions; adopting its name does not
transfer every result to a changed encoding.

There is a plausible repair direction, which should be distinguished
from a demonstrated mode failure: with an exact, non-wrapping bit count,
a correctly authenticated final flag, canonical partial bytes, and a
defined empty-message case, the written message-to-transcript encoding
is injective. Unequal total lengths differ in the final count;
equal-length distinct messages differ in their data. Zero padding is
therefore not intrinsically fatal to the written design. R1 is the
concrete case where an implementation removes an essential part of that
argument.

Second, even within its stated ideal model, the simulator description is
incomplete. Its accounting for unreachable queries, later connections
into the reachable graph, and out-of-order queries is not sufficient to
justify the claimed perfect consistency. The stated B2 event also needs
to distinguish the public IV from guessed unknown states. The query
measure must account for the total number of compression evaluations or
processed blocks: one variable-length hash query can process many
blocks. A complete theorem needs an exact interface, game, simulator,
resource measure, and bad-event proof.

Third, the displayed advantage bound \(O(q^2 / 2^{1024})\) becomes
non-informative around the 512-bit birthday scale. Even a correct
version of that bound would not, by itself, establish a 1,024-bit
preimage or second-preimage claim. Those properties need their own
bounds and input-length/resource assumptions. Conversely, the
birthday-scale limit of this proof is not itself a preimage attack;
these are different statements.

Fourth, §4.4 changes idealizations. The
\href{https://www.cs.ucdavis.edu/~rogaway/papers/hash.pdf}{Black--Rogaway--Shrimpton
analysis} gives black-box results for block-cipher-based hashing under
its stated model. It does not prove that FEILIAN's repeated-message ARX
structure is an ideal cipher, or automatically identify Davies--Meyer
compression with the independent random compression oracle assumed by
§4.3.1. The acknowledged Davies--Meyer fixed-point property reinforces
the need to keep those interfaces and assumptions distinct. For a freely
chosen compression state, its cost is comparable to one inverse of the
pre-feedforward permutation, not a \(2^{512}\) birthday search. Reaching
such a state from the fixed hash IV under valid counters, or
establishing a fixed point of the separate XOF mode, are different
questions. No such reachability result or witness is established here.
No complete composition argument is provided.

Fifth, the component evidence does not quantify a full-compression
security margin. The statistical tests in §4.2 and average single-bit
avalanche table in §4.5.3 measure useful implementation properties, but
cannot certify collision or preimage complexity. They do not establish
bounds on worst-case structured differences, differential hulls, related
message/tweak behavior, or full-round attacks. NIST itself cautions
against using its
\href{https://www.nist.gov/news-events/news/2022/04/decision-revise-nist-sp-800-22-rev-1a}{SP
800-22 statistical suite as a cryptographic security assessment}.

The critical composition questions are:

\begin{itemize}
\tightlist
\item
  The same message is injected five times, with no expansion or word
  reordering. Security analysis must include those correlated inputs;
  analyzing message-free permutations expressly leaves out that freedom.
\item
  Constants repeat within each phase and the first and fifth phase have
  the same constant pattern. The phases are not independent random
  permutations.
\item
  Column-wise SubColumn and row shifts retain a structural relationship
  with cyclic column permutations. Rotating the constants between phases
  makes joint analysis of the message and tweak essential. Whether a
  related parameter setting is a valid hash invocation depends on
  version, flag, length, and counter-word conventions.
\item
  Reaching an attacker-selected internal state from the fixed IV is a
  real constraint, but the assertion that this necessarily costs at
  least a birthday search is not a general theorem. Likewise, secrecy of
  a chaining state alone does not prove every keyed application secure.
\item
  The final message injection occurs in round 17, leaving seven
  SubColumn/ShiftRow layers after that injection. That is not evidence
  of an attack, but it explains why counting all twenty rounds is an
  incomplete way to state the margin against message-dependent methods.
\end{itemize}

The discussion of slide attacks in §4.6.2 correctly recognizes repeated
structure but then appeals to diffusion and nonlinear operations making
a slid relation improbable. Those properties alone do not exclude
structural slide relations. The meet-in-the-middle and boomerang
discussions similarly assert adequate phase counts without supplying
quantitative bounds or reproducible supporting searches. These are
hypotheses to investigate, not established exclusions.

There are also checkable errors in the security text. On p.~25 the
displayed boomerang criterion is \(p^2q^2 > 2^{2n}\), which is
impossible for probabilities. This was visually verified in the PDF.
Correcting the missing minus sign would not resolve the more general
issue: the appropriate random baseline and the use of
\texttt{p\^{}2\ q\^{}2} require a precisely defined experiment and
justification of dependencies. Reference {[}9{]}, cited for a hash
second-preimage argument on p.~20, is a paper about KeeLoq rather than
the claimed hash result. The unsupported numerical claims and
mis-citation should be corrected, not used as security evidence.

There are also positive mathematical checks. Both Sigma maps have rank
64 over GF(2). Representing rotation by \texttt{x}, the word modulus is
\(x^{64}+1=(x+1)^{64}\); each three-term Sigma polynomial takes value 1
at \texttt{x\ =\ 1}, so it is coprime to that modulus and invertible.
Sage independently confirmed the gcds and one-dimensional fixed spaces.
These facts do not rely on one rotation amount being prime. They
establish linear-map properties, not differential probabilities through
modular additions.

The SubColumn inverse in §4.5.2 is algebraically consistent with the
forward equations, and 10,000 deterministic random 256-bit inputs passed
a forward/inverse check. The algebraic inverse, together with invertible
row shifts and XOR injections, establishes that the round composition
before feedforward is a permutation for fixed message and tweak. There
is no rank defect in either Sigma map. The probability-one component
differentials acknowledged by the specification remain relevant to
security analysis, but their existence alone does not establish a
full-hash weakness. No independent exhaustive differential search or
full-round attack construction was performed here.

The generic-security discussion should distinguish the variants. Only
the 1,024-bit output version has a chaining state equal in width to the
digest. The 512- and 768-bit versions retain a 1,024-bit state
internally. A collision in a truncated output is not necessarily an
internal-state collision, so generic multicollision and herding costs
cannot be transferred between these cases merely by substituting the
digest length. The fixed-salt HAIFA analysis also retains qualifications
concerning multicollisions and herding; an unqualified claim of
inheriting all HAIFA protections is too broad.

\Needspace{5\baselineskip}

\subsection{Classical and quantum
claims}\label{classical-and-quantum-claims}

Table 4.1 explicitly labels its claims as classical. Under ideal
behavior, its collision and preimage exponents are the usual ones for
the three output lengths; their achievement by FEILIAN is what remains
unestablished. The generic quantum query scales are different:

\Needspace{7\baselineskip}

{\def\LTcaptype{none} % do not increment counter
\begin{longtable}[]{@{}
  >{\raggedleft\arraybackslash}p{(\linewidth - 8\tabcolsep) * \real{0.1200}}
  >{\raggedleft\arraybackslash}p{(\linewidth - 8\tabcolsep) * \real{0.2200}}
  >{\raggedleft\arraybackslash}p{(\linewidth - 8\tabcolsep) * \real{0.2200}}
  >{\raggedleft\arraybackslash}p{(\linewidth - 8\tabcolsep) * \real{0.2200}}
  >{\raggedleft\arraybackslash}p{(\linewidth - 8\tabcolsep) * \real{0.2200}}@{}}
\toprule\noalign{}
\begin{minipage}[b]{\linewidth}\raggedleft
Digest bits
\end{minipage} & \begin{minipage}[b]{\linewidth}\raggedleft
Classical collision exponent
\end{minipage} & \begin{minipage}[b]{\linewidth}\raggedleft
Classical preimage exponent
\end{minipage} & \begin{minipage}[b]{\linewidth}\raggedleft
Generic quantum collision exponent
\end{minipage} & \begin{minipage}[b]{\linewidth}\raggedleft
Generic quantum preimage exponent
\end{minipage} \\
\midrule\noalign{}
\endhead
\bottomrule\noalign{}
\endlastfoot
512 & 256 & 512 & about 170.7 & 256 \\
768 & 384 & 768 & 256 & 384 \\
1,024 & 512 & 1,024 & about 341.3 & 512 \\
\end{longtable}
}

The quantum columns use the usual idealized collision-search and
preimage-search query models, from
\href{https://arxiv.org/abs/quant-ph/9705002}{Brassard--Høyer--Tapp} and
\href{https://arxiv.org/abs/quant-ph/9605043}{Grover}. They are not
practical gate-count or memory estimates and are not newly discovered
attacks on FEILIAN. The specification's claim that there is no better
dedicated quantum method needs analysis of the actual construction; a
classical simulator proof does not establish it. In particular, the
512/1,024-bit classical claim should not be restated as that many bits
of quantum security.

\Needspace{5\baselineskip}

\section{R7. MAC, KDF, XOF and interface
composition}\label{r7.-mac-kdf-xof-and-interface-composition}

The claim that the preceding indifferentiability discussion licenses
general replacement of a random oracle is too broad even before the
proof defects are considered.
\href{https://eprint.iacr.org/2011/339.pdf}{Ristenpart, Shacham, and
Shrimpton} show why ordinary indifferentiability does not automatically
cover all security notions with multiple adversarial stages.

For the proposed prefix MAC, finalization is a sensible defense against
ordinary length extension, but collision resistance alone does not prove
MAC security. A useful bound must include key entropy, tag length, query
counts, total processed data, and the allowed interfaces. A claim of
512-bit MAC security cannot hold independently of those parameters. A
vetted MAC construction would still require a sound assumption about
FEILIAN and a correct implementation.

Tag collisions among already queried messages are not fresh-message MAC
forgeries. For an ideal random-function MAC with \(t\)-bit tags, the
generic verification-guessing contribution is bounded by \(q_v/2^t\),
where \(q_v\) counts verification attempts. A PRF-based argument adds
the relevant distinguishing advantage; see
\href{https://www.cs.ucdavis.edu/~rogaway/classes/227/spring05/book/main.pdf\#page=167}{Bellare--Rogaway,
Proposition 7.3}. The \(2^{t/2}\) birthday scale for tag collisions does
not, by itself, reduce the generic forgery exponent to \(t/2\). This
distinction does not supply the missing proof for FEILIAN's actual
prefix MAC.

For the KDF, \texttt{PRK\ =\ FEILIAN(IKM)} does not establish a general
extractor for arbitrary non-uniform sources. A deterministic function
cannot guarantee uniformly distributed output for every high-entropy
source; a random-oracle argument requires suitable independence,
secrecy, and entropy conditions. Hashing does not create entropy. The
proposed expansion also needs an unambiguous encoding of context and
counter, an output limit, and explicit separation from other uses of the
same secret. A fixed-length PRK and a fixed, known-width trailing
counter give unambiguous parsing even when \texttt{info} is
variable-length: the counter is the final fixed-width field. Specify its
injective encoding, byte order, allowed iteration range starting at one
and overflow rejection. A variable-width alternative needs its own
unambiguous framing. The specification gap does not imply that every
possible encoding is ambiguous. Calling PRK a salt is misleading: in
that expression it serves as secret key material. This is not HKDF,
whose extract step is HMAC-based and whose salt and context handling are
specified in \href{https://www.rfc-editor.org/rfc/rfc5869.html}{RFC
5869}.

For the XOF, the distinct squeezing version values are a useful design
choice. However, the mode changes the finalization rule during
absorption and then applies compression to an internal root; it is not
the hash mode analyzed in §4.3.1. It needs its own definition and proof.
Specify the empty-block encoding, counter word order and bound, the
output truncation for each variant, the initial root for an empty input,
and behavior on counter exhaustion. A finite 128-bit counter cannot
silently be allowed to wrap while retaining the claim that every squeeze
call is distinct. No MAC/KDF/XOF implementation or corresponding KATs
were present.

Keyed applications create additional implementation requirements. In
1SC, 2SC, and 4SC, the digest signal is continuously derived from the
current chaining register, and the top-level MMIO read path is not
restricted to completed hashes. Message registers are readable too.
Whether this crosses a security boundary depends on who can access the
peripheral; it is not an additional public-hash vulnerability by itself.
It does mean the supplied public-hash interface must not be assumed to
keep a secret prefix or intermediate keyed state isolated. Similarly, C
heap copies of keyed inputs are freed without explicit erasure. A keyed
API should define state exposure and clearing rather than inherit them
accidentally.

The scalar and SIMD ARX arithmetic has fixed operation schedules and no
obvious message-dependent branches or memory indices. That is favorable
source-level evidence for timing behavior on the intended platforms. No
machine-code timing certification, leakage measurement, fault analysis,
or secret-state isolation validation was performed, so these are not
established properties of a deployable MAC or KDF.

\Needspace{5\baselineskip}

\section{R8. Shared DRNG and test-driver
assurance}\label{r8.-shared-drng-and-test-driver-assurance}

The six \texttt{drng.c} files supplied with FEILIAN are byte-identical,
as recorded in \href{data/shared_drng.json}{shared\_drng.json}. Its
header attributes the software to ICCS, and the
\href{data/SOURCES.md\#source-15}{submission README} marks the KAT
generator and DRNG files as unmodifiable submission infrastructure.
These observations support treating the defects below as a shared
test-infrastructure issue and coordinating an upstream correction with
the organizers. They are not defects in FEILIAN's compression
mathematics.

The \href{data/SOURCES.md\#source-10}{SM3 rotation macro}, at line 36,
allows a shift by the word width when the rotation count is zero. The
SM3 round calls at lines 90 and 94 reach that case. With 32-bit
\texttt{unsigned\ int}, this is undefined behavior under C's shift
rules, regardless of whether a particular processor masks its
machine-instruction shift count. Use defined 32-bit rotation semantics
for every permitted count; see
\href{https://www.open-std.org/jtc1/sc22/wg14/www/docs/n1570.pdf}{C11
draft N1570, §6.5.7 paragraph 3}.

Error propagation is also incomplete. \texttt{SM3\_DRNG\_Instantiate}
ignores both \texttt{SM3\_df} return values at lines 252 and 257, and
the \href{data/SOURCES.md\#source-14}{KAT driver} ignores the
initialization status at lines 296, 379, 487 and 601. The failed state
is \textbf{not necessarily all zero}. Failure of the outer allocation
follows a context clear; an inner derivation failure can leave seed
material or a padded state copied into the context before success is
returned. The result depends on the failing stage and input. Propagate
every initialization/derivation failure and stop generation before
recording purported test vectors.

The partial-byte mask used by the generator is consistent with its
retained-bit count; it is not an additional mask defect. The accepted
seed-pointer contract should also be explicit: if a null pointer with
zero length is supported, avoid passing it to \texttt{memcpy}; C's
library pointer requirements still apply at zero length (§7.24.1
paragraph 2). The supplied normal KAT initialization passes a real seed
array, so this is conditional API hardening, not a demonstrated
normal-input failure.

These conclusions are static. No sanitizer run or allocation-failure
experiment from the second review is claimed as independently
reproduced. They affect the portability and reliability of regenerating
vectors. They do \textbf{not} invalidate the already supplied
message/digest pairs or the 24,582 comparisons performed here, which
hash fixed KAT message bytes directly and do not execute the DRNG.
Preserve those records while correcting and independently checking the
generator.

\Needspace{5\baselineskip}

\section{Validation and limits}\label{validation-and-limits}

The \href{run.py}{verification driver} can build the six supplied C
sources from a separately obtained, hash-verified submission and checks
canonical inputs only. The \href{evidence/c_conformance.json}{C
conformance record} and \href{evidence/quick.json}{model record} contain
the results and complete ordinary-example digests. The
\href{code/check_appendix.py}{appendix check} verifies the six printed
examples under their identified count conventions, all 45 round
snapshots, nine initial states and nine recovered domains. The
\href{code/arithmetic_checks.sage}{Sage checks} perform exact polynomial
and constant calculations, with a separate
\href{evidence/arithmetic.json}{arithmetic record}. The
\href{data/rtl_source_facts.json}{RTL source index} identifies the
relevant counter and variant locations; \href{data/SOURCES.md}{source
locators} also cover the hardware interface and shared test generator.
Their source observations are not executable hardware or failure-path
tests.

The package provides a self-contained Python check of ordinary examples
and component arithmetic, plus an optional check against the supplied C
sources and published short-message KATs. See \url{REPRODUCING.md} for
commands, dependencies and coverage. The latter path verifies all 74
original file hashes before compilation. Original submission files are
not redistributed.

The ordinary C tests passed. Compiler warnings concerned a validated
positive output-length conversion and unused SIMD masks; those warnings
were not treated as security findings. The long-message KAT sets
(\(2^{23}\), \(2^{33}\)) and million-iteration loop were not run. No
software source was patched, no noncanonical bit-input or failure-path
witness was exercised, and no HDL execution was performed. Passing these
checks is evidence about specified computations and internal
consistency, not a proof of security.

\Needspace{5\baselineskip}

\section{Requested corrections}\label{requested-corrections}

Request a corrected submission before spending substantial effort on
full-round security claims. The immediate requests should be:

\begin{enumerate}
\def\labelenumi{\arabic{enumi}.}
\tightlist
\item
  Designate one authoritative algorithm, settling the IV, AddConstant
  rows, state/byte order, counter word order, empty input, partial-bit
  convention, and all named output variants. Identify which exact
  definition produced each security experiment.
\item
  Correct the RTL length binding and variant definitions, and provide
  end-to-end tests across all four architectures against an
  independently implemented model. Compression-only synthesis results do
  not validate message framing or digest presentation.
\item
  Correct the C input-length checks, partial-byte handling, allocation
  error propagation, and output-length contract. Prefer a bounded-memory
  streaming interface with explicit counter-overflow behavior.
\item
  Regenerate all appendix and KAT material from the designated version
  using true message-length binding, with independently checked
  intermediate states and a reproducible provenance record. Explain the
  mixed appendix counter conventions.
\item
  Coordinate correction of the shared DRNG and KAT driver with the
  organizers: use defined rotations, propagate failures and
  independently check deterministic output across supported toolchains.
  Preserve existing vectors for comparison.
\item
  Replace the mode proof with one for the actual interface and obtain
  independent review of the complete message/tweak/state composition.
  State separately what is proved in an ideal model, what is
  experimentally supported, and what remains conjectural. MAC, KDF, and
  XOF recommendations need their own specifications and bounds.
\end{enumerate}

The hardware encoding defect is the strongest concrete reason to stop
use of the affected implementations now. The inconsistent algorithm
definitions and unsupported security margins are the strongest reasons
to withhold confidence in the submission as a whole. They justify a hold
and a corrected review target; they do not justify claiming that this
review has broken the full-round, correctly encoded software hash.

\Needspace{5\baselineskip}

\section{References}\label{references}

\begin{enumerate}
\def\labelenumi{\arabic{enumi}.}
\item
  Lei Wang, Ling Song, Yaobin Shen, Kaixuan Wang, Zicheng Shi and Xin
  Yi. \emph{FEILIAN algorithm specification}, 30 June 2026. The supplied
  68-page PDF is identified by its digest above and in
  \href{data/input_manifest.json}{the input manifest}. Printed page
  numbering is used throughout.
\item
  Markku-Juhani O. Saarinen. \emph{FEILIAN (hash-10)}, 23 September
  2026. \href{https://ngcc.dev/reports/hash-10.html}{Public report},
  consulted 24 September 2026.
\item
  Eli Biham and Orr Dunkelman. \emph{A Framework for Iterative Hash
  Functions --- HAIFA}. IACR ePrint 2007/278.
  \href{https://eprint.iacr.org/2007/278.pdf}{Paper}.
\item
  John Black, Phillip Rogaway and Thomas Shrimpton. \emph{Black-Box
  Analysis of the Block-Cipher-Based Hash-Function Constructions from
  PGV}. CRYPTO 2002.
  \href{https://www.cs.ucdavis.edu/~rogaway/papers/hash.pdf}{Author-hosted
  paper}.
\item
  NIST. \emph{Decision to Revise NIST SP 800-22 Rev.~1a}, 19 April 2022.
  \href{https://www.nist.gov/news-events/news/2022/04/decision-revise-nist-sp-800-22-rev-1a}{Notice}.
\item
  Lov K. Grover. \emph{A Fast Quantum Mechanical Algorithm for Database
  Search}. 1996. \href{https://arxiv.org/abs/quant-ph/9605043}{Paper}.
\item
  Gilles Brassard, Peter Høyer and Alain Tapp. \emph{Quantum
  Cryptanalysis of Hash and Claw-Free Functions}. 1997.
  \href{https://arxiv.org/abs/quant-ph/9705002}{Paper}.
\item
  Thomas Ristenpart, Hovav Shacham and Thomas Shrimpton. \emph{Careful
  with Composition: Limitations of the Indifferentiability Framework}.
  EUROCRYPT 2011. \href{https://eprint.iacr.org/2011/339.pdf}{Paper}.
\item
  Hugo Krawczyk and Pasi Eronen. \emph{HMAC-based Extract-and-Expand Key
  Derivation Function (HKDF)}. RFC 5869, May 2010.
  \href{https://www.rfc-editor.org/rfc/rfc5869.html}{RFC}.
\item
  ISO/IEC JTC1/SC22/WG14. \emph{N1570: Committee Draft, Programming
  Languages --- C}, 12 April 2011. §§6.5.7 and 7.24.1.
  \href{https://www.open-std.org/jtc1/sc22/wg14/www/docs/n1570.pdf}{Draft}.
\item
  Mihir Bellare and Phillip Rogaway. \emph{Introduction to Modern
  Cryptography}, course notes, 2005. Proposition 7.3.
  \href{https://www.cs.ucdavis.edu/~rogaway/classes/227/spring05/book/main.pdf\#page=167}{Author-hosted
  notes}.
\end{enumerate}

The original text of this report is licensed under CC BY 4.0. See
\url{LICENSING.md} for software terms and third-party exceptions.

\end{document}
